Now I have enough information to make the necessary corrections. Let me analyze the key issues:

1. **CLT vs LLN**: Chapter 4 explicitly distinguishes CLT (about distribution shape) from 大数の法則 (about convergence to population mean). The text in 9.2 attributes "標本平均は母平均に近づく" to CLT, but this is the LLN statement.

2. **t-statistic formula**: Chapter 6 defines the t-statistic as the two-sample form $(\bar{x}_1 - \bar{x}_2)/\mathrm{SE}$, but chapter 9 references "$t = \bar{x}/\mathrm{SE}$" which is a different (one-sample, μ_0=0) form.

3. **Use of "独立" for confounding**: In a section that has just discussed statistical independence, calling a confounder's effect "独立に押し上げている" risks confusion with statistical independence.

Now let me output the corrected text:

\section{統計の限界・バイアス・倫理：手法の外側を知る}

\subsection{1章で告げた「正しく計算しても間違う」の正体}

道具は揃った.
2章で平均と散らばりを言葉にし, 5章で幅をつけ, 8章で関係を直線で表す. これでだいたいの分析はできる.

ただし, ここまでの道具を全部正しく使っても, 結論が間違うことがある.
1章のはじめで, そういう例があるとだけ告げてあった.
あれの答え合わせを, ここでする.

1973年, カリフォルニア大学バークレー校の大学院は, 性別を理由に女子受験生を差別しているとして訴えられた.
証拠とされたのは合格率の差だ.
全志願者を集計すると, 男子の合格率は約44\%, 女子は約35\%.
9ポイントの差は無視できない.

この数字だけ見れば, 大学院に問題があると考えるのが自然だ.
ところが大学院側が学部別に集計し直したところ, 6つの主要学部のうち4つで, 女子の合格率は男子と同じか, むしろ高かった.
残り2つの学部でも数ポイント差で男子のほうがわずかに高い程度で, どこをどう見ても「女子だから不利」と言える状況ではなかった.

なぜそうなったのか. 応募の分布が学部によってまるで違ったのだ.
女子は合格率がそもそも低い学部（人文系で20\%から30\%）に多く志望し, 男子は合格率が高い学部（理工系で60\%から70\%）に多く志望していた.
各学部の中では公平でも, 女子が「狭き門の学部」に集中していたために, 全体平均では女子の合格率が下に引かれた.

集計の単位を変えただけで結論が逆転する.
このような現象を\textbf{シンプソンのパラドックス}と呼ぶ.
平均は2章で, 散布や相関は7章で扱った. どの計算も正しい.
なのに, 全体集計と層別集計で正反対のことが言える.

この種の食い違いは珍しい例外ではない.
第3変数が群によって違う割合で混ざっていれば, 同じ仕掛けはどんな分析でも起こりうる.
バークレーの例なら「学部」が第3変数で, 男女で志望先の分布が違っていた.
シンプソンの匂いを嗅ぐ手がかりはわりとはっきりしている. 全体での集計と, 重要な属性で層別した集計を必ず両方確かめることだ.
両者が食い違っていれば, 全体集計のほうが第3変数で歪んでいる可能性が高い.

8章の最後で,「手法の外側を知ることが, 手法を正しく使うための最後の一歩になる」と書いた.
手法の外側とは, おおまかに言って3つの方角だ.
標本をどう集めたか. 集まったデータの中で何が見えていないか. 出てきた数字をどう解釈したか.
本章ではこの3方向を順に歩き, 続いて手法そのものが暗黙に置いている前提に踏み込む.
最後に, 数字を扱う側の責任の話で本書を閉じる.

まずは標本の集め方から始めよう.
3つのうち最も入口に近く, それでいて最も誤解されやすい場所だからだ.


\subsection{バイアスの三つの顔：標本・観測・解釈}

データに歪みが入る場面は, 大きく分けて3つある.
標本を集める段階, 集めたデータの中身が偏っている段階, 集まったデータを解釈する段階.
データの流れの上流・中流・下流に, それぞれ違う種類の歪みが棲んでいる.

\subsubsection*{上流：サンプリングバイアス}

5章で信頼区間を計算したとき,「件数を増やせば幅は狭くなる」という結果を手に入れた.
$n$を増やせば$\sqrt{n}$で精度が上がる. 計算式の上ではその通りだ.
ただし, この話は標本が母集団から無作為に抽出されているという前提の下でだけ成り立つ.

抽出が偏っていたらどうなるか.
たとえば, 知りたいのは「日本に住む有権者」の意見なのに, 集めたのは「平日の昼に渋谷駅にいた人」だとしよう.
件数を10万人に増やしても, 平日の昼に渋谷駅にいない人の意見はゼロのままだ.
標本平均が近づいていく先は,「平日の昼に渋谷駅にいた人たち」の真の平均でしかない.
4章の大数の法則で「標本平均は母平均に近づく」と言ったとき, その母平均は標本を取った母集団の平均であって, 知りたかった母集団のそれとは限らない.

これが\textbf{サンプリングバイアス}である. 件数で殴って解決できる問題ではない.

5章で紹介したリテラリー・ダイジェスト誌の話を思い出してほしい.
1936年の大統領選で, 同誌は約240万人を対象とした大規模調査からランドンの勝利を予測した.
ところが結果はルーズベルトの圧勝だった. 同じ年, ジョージ・ギャラップは5万人の調査で正しく予測している.
件数では数十倍劣るギャラップが勝ったのは, 標本の代表性で勝っていたからだ.
ダイジェスト誌のリストは電話帳と自動車登録名簿から作られていて, 大恐慌下で電話も車も持たなかった貧困層をまるごと外していた.
件数の多さは偏りの方向に対して何もしてくれない. この事例はそれを教えている.

\subsubsection*{中流：生存者バイアス}

標本の代表性は,「取れたデータの中身」を眺めても判断できないことがある.
そもそも取れなかったデータがあると気づかないと, 結論は静かに歪む.

第二次大戦中, 米軍は帰還した爆撃機の被弾箇所を調べて, どこを補強すべきかを検討していた.
集計してみると, 翼の先端と機体後部に弾痕が集中している.
最初の判断は素直だった.「弾痕が多い場所を補強しよう」.

統計学者アブラハム・ワルドはこの結論を止めた.
「補強すべきは, 弾痕の少ない場所のほうだ」.
帰還した機は, 翼や機体後部に被弾しても飛び続けられた機だ.
エンジンや操縦席にやられた機は, そもそも基地に戻ってこない.
分析の対象に上がっていたのは「生き残った機」だけで, 撃墜された機の被弾箇所はデータに現れていなかった.

これが\textbf{生存者バイアス}である.
「観測できたものだけ」を見て結論を出すと, 観測できなかったものが体系的に抜け落ちる.
事業継続中の企業だけを集めて成功要因を語る研究, 在籍中の社員だけを集めた満足度調査, 完走したマラソン選手だけを集めた走法分析.
どれも, 退場した側に重要な情報が眠っている可能性がある.

データの中身を眺めるだけでは, 何が抜けているかは見えない.
「このデータには本当はどんな対象が含まれるべきか, そのうち実際に観測できたのは誰か」を, データを見るのとは別の頭で考える必要がある.

\subsubsection*{下流：確証バイアス}

仮にデータの集め方が完璧で, 観測の漏れもなかったとしよう.
それでも, 出てきた数字をどう解釈するかの段階で歪みが入る.

自分が正しいと思っている仮説に合うデータは,「やはり」と受け取ってしまう.
合わないデータは「外れ値だろう」「サンプルが悪かったのだろう」と片付けたくなる.
こういう経験は誰にでもある. これに名前をつけたのが\textbf{確証バイアス}だ.

統計の道具で正当化された確証バイアスもある.

公平なコインを20枚用意して, それぞれ100回投げ, 各コインで「表と裏の出方が偏っているか」を有意水準5\%で検定すると考えてみよう.
20枚すべてが本当に公平だったとしても, 偶然どれか1枚は「有意」と出る.
1枚あたり有意にならない確率が0.95, 20枚すべてで有意にならない確率が$0.95^{20} \approx 0.36$, したがってどこか1枚以上で偽陽性が出る確率は$1 - 0.36 \approx 0.64$.
20回試せば, 6割以上の確率で「有意な発見」がでっち上がる.

ここで実験者が「有意になった1枚」だけを論文にし, 残り19枚を引き出しにしまったらどうなるか.
表に出るのは「特殊なコインを発見した」という主張だけだ.
6章で習ったp値の計算は, どの1枚についても正しく行われている.
それでも結論は嘘になる.

これが\textbf{p-ハッキング}と呼ばれる行為である.
6章で覚えたp値が, ここでは正当化の道具として悪用されている.
都合のいい結果が出るまで分析の角度や対象を変え続けると, 確証バイアスは自動化される.

防ぎ方の入口だけ書いておく.
分析を始める前に「どの仮説をどう検証するか」を確定させて記録する事前登録（pre-registration）, 複数の検定をまとめて行ったときに有意水準を厳しくする多重比較の補正, データと手順を公開して第三者の再分析に晒すこと.
どれも本書の範囲を超えるが, 道はこの方向に開いている.

ここまで3つの歪みを並べた.
標本の集め方（上流）, 観測対象の選別（中流）, 解釈の仕方（下流）.
データが流れる順番に沿って, 別々の場所に違う罠が仕掛けられている.

3類型は別々の知識ではなく, 1本の流れの上の地点だ.
分析を始める前にまず上流を疑い, 手元のデータを見ながら中流を疑い, 結果を読むときに下流を疑う.
本章の残りはこの順序のまま進む.
次は, 集まったデータに手法を当てる段階, つまり「中流から下流の入口」に位置する話に踏み込む.


\subsection{前提が崩れるとき：独立性・線形性・等分散性}

データの集め方や選別の歪みを乗り越えたとしよう.
それでも, 手法そのものが暗黙に置いている前提が崩れていれば, 出てきた数字の意味は変わる.
2章から8章で使ってきた手法のうち, 推論に関わる部分は概ね3つの前提に乗っている.
データ点が独立であること. 関係がおおむね直線で表せること. 残差の散らばりがどの場所でも同じくらいであること.

\subsubsection*{独立性：隣の点と関係していないか}

4章のCLTの前提として「独立同分布」を導入した.
ここでの「独立」は, データ点どうしが互いに無関係であることを指す.
ある月のデータが翌月のデータに影響していたら, 独立ではない.

8章で扱った広告費と売上の例を思い出してほしい.
あの10ヶ月分のデータは時系列に並んでいた.
回帰の計算式は, 各データ点を独立に扱って残差の二乗和を最小化している.
ところが, 先月の広告効果が今月にも残っていたとしたらどうなるか.
ある月の売上が高ければ, その勢いで翌月も高い. 残差の符号は隣どうしで似てしまう.

このとき計算式は止まらない. 傾きも切片も決定係数も, 数字としては出る.
ただし, 出てきたSEは「独立だった場合に成り立つ式」で計算されているから, 独立でない実態に対しては不当に小さい値になる.
SEが小さいと, 信頼区間も狭く出る.
p値はどう動くか. 6章のt統計量は$(\bar{x}_1-\bar{x}_2)/\mathrm{SE}$型の比であり, SEが分母に入っている.
分母が小さくなれば$t$は大きくなる. $t$が大きい側に振れれば, それに対応するp値は小さく出る.
本当は不確かな結論が, 紙の上では確実そうに見えるわけだ.

時系列データだけではない. クラスごとの生徒, 同じ会社の社員, 同じ地域の住民. 集団内の点は互いに似ている傾向がある.
集団効果を無視して全員を独立だと扱うと, 同じ罠が起こる.
独立性は, 計算が動くかどうかではなく, 計算結果の信頼性を決めている.

崩れていた場合の道筋だけ名前を出しておく. 時系列の依存があるなら自己回帰モデル, 集団内で似ているならクラスタロバスト標準誤差や階層モデル. 8章までの直線回帰の延長で扱える話だ.

\subsubsection*{線形性：直線で表せる関係なのか}

8章で残差プロットを導入したとき, アンスコムの第2組の話をした.
4組のデータが$r \approx 0.82$, $R^2 \approx 0.67$という似た数字を返したが, 第2組は実際にはきれいな放物線で, 直線を当てはめること自体が無意味だった.

本当の関係が放物線, 指数, 周期的なものだったら, 直線を引いても得られるのは「もっとも近い直線」だけだ.
傾きや切片の数値は出る. しかしそれらは現実の関係を要約していない.

線形性が崩れているかどうかは, 残差プロットで目視できる.
予測値の小さい所と大きい所の両方で残差が正に偏り, 中央で負に偏っていれば, 関係はU字型だ.
逆向きのパターンも同様の警告だ.
8章で「ゼロのまわりにランダムに散らばっていればよい」と書いたのは, この前提が満たされているかを目で確かめる手順だったわけだ.

線形性の確認は, 数値だけ見ていても気づけない.
$R^2$が0.7を超えていても, 残差プロットが系統的なパターンを示していれば, 直線モデルは不適切だ.
逆に$R^2$が0.4でも, 残差にパターンがなければ, それはそれで意味のある直線関係である.
$R^2$は当てはまりの「総量」しか語らず, 「外し方の質」は残差プロットでしかわからない.

非線形が見えたら, $\log x$や$x^2$といった変数変換を入れた回帰や, 多項式回帰・非線形回帰に進む手がある. 直線という発想を完全に手放すわけではなく, 直線が当てはまる空間を探しに行く方向だ.

\subsubsection*{等分散性：場所によって散らばりが違わないか}

3つめの前提が, 残差の散らばりがデータの位置によって変わらないこと, つまり等分散性である.
6章のt検定でWelch法を採ったのは, 2群間で分散が違うことを許す扱いを選んだという意味で, この前提を緩めた選択だった.

回帰では事情が少し違う.
広告費が大きい月ほど売上の上下幅も大きくなるとき, 残差プロットには左端が密で右に行くほど縦に広がる扇形が現れる.
直線そのものは間違っていない. 平均的な傾きは正しく推定されている.
ただ,「来月の広告費を50万円にしたら売上はおよそいくつか」と幅つきで予測したい場合, 一律のSEで予測区間を出すと, 場所によって狭すぎたり広すぎたりする.
平均値の推定には使えても, 予測の幅は当てにならなくなる.

不均一な散らばりが見えたときは, 場所ごとに重みを変える重み付き回帰や, 不均一性に頑健なロバスト標準誤差で受け止める方法がある.

ここまでの3つの前提を点検する手段は, ほとんど残差プロット1枚に集約される.
残差を予測値に対してプロットし, ランダムに散らばっていればよし, 系統パターンや扇形が見えたら警告.
8章で導入したあのプロットは, 単なる確認作業ではなく, 線形回帰という道具を使ってよいかどうかを判断する診断装置だったのだ.

ここまでで, データの集め方と手法の前提は点検する目を持った.
それでもまだ,「正しく計算しても結論が間違う」場面がある.
計算自体は正しく動き, 前提も満たされている. それなのに, 結果の読み方が間違っているケースだ.
次節はその話に移る.


\subsection{正しく計算しても結論が崩れるとき：p値・相関・外挿}

ここまでは, データの上流から手法の前提までの話だった.
すべての段階を正しく済ませたとしよう.
それでも, 出てきた数字の解釈が間違っていれば, 結論はやはり間違う.
誤読が起きやすい場面を, p値, 相関と因果の関係, データ範囲の外側の順に見ていく.

\subsubsection*{p値の周りで誤読が起きる}

6章でp値を「帰無仮説が正しいと仮定したとき, 観測値以上に極端な結果が出る確率」と定義した.
この定義はそれなりに込み入っているせいで, 短く要約しようとすると意味がずれる.
よくある誤読を, 出会う頻度の高い順に並べておく.

おそらく最も多いのは,「p値0.03は帰無仮説が間違っている確率97\%である」と読む誤りだ.
これは違う. p値は「帰無仮説が正しいと仮定したとき」の珍しさを測っているのであって, 帰無仮説そのものの真偽の確率を返してくれるわけではない.
p値が0.03でも, 仮説が間違っている確率は3\%ではない.

次に多いのが,「p<0.05なら効果は大きい」と読む誤りだ.
p値は珍しさの指標であって, 効果の大きさは別の指標で測る.
2群の平均差を比べる場面で, 群1の平均が3.04, 群2が3.01, 差は0.03点だったとしよう.
各群の標本数が10000あれば, この差はp<0.05で有意になる.
件数が大きいと, 実用上は無視できる差でも「珍しい」と判定される.
逆に件数が少なければ, 大きい差でも有意にならないことがある.

最後に,「p>0.05なら効果はない」と読む誤り.
これも違う.「珍しいと言える証拠が足りない」と「効果がないことが証明された」は別の話だ.
n=20で実験して有意にならなかったからといって, 効果が存在しないわけではない.

p値の話を扱うときは, p値だけで報告を済ませない方がよい.
5章で覚えた信頼区間と, 6章Part Bで触れた効果量を一緒に出す.
信頼区間は「効果がだいたいどの範囲にありそうか」, 効果量は「差そのものの大きさ」, p値は「偶然で説明しにくいかどうかの目安」.
それぞれ別の問いに答えている.

報告の側もこれをセットで書き, 読み手の側もこれを求めるとよい.
p値だけが踊っている結論は, それだけで疑う材料になる.

\subsubsection*{相関の正体は何種類かある}

7章と8章で「相関は因果を意味しない」と何度か言った.
言葉としては理解しやすい. けれど, 実際に強い相関を見たとき, それが因果でないなら何なのかは, きちんと整理しておく価値がある.

第3変数が裏で両方を動かしているケースから始めよう. 名前を\textbf{交絡}という.
アイスクリームの月別消費量と水難事故の月別件数には, 強い正の相関がある.
アイスを食べると溺れるわけではない. 気温が高い月にアイスもよく売れ, 同じく気温が高い月に水場へ行く人が増えて事故も増える.
両方を別々に押し上げている第3変数（気温）があるとき, アイスと水難事故の間に直接の因果がなくても相関は出る.
この第3変数を交絡因子と呼ぶ.

ヨーロッパ各国でコウノトリの巣の数と人間の出生数に正の相関がある, という有名な話もこの仲間だ.
コウノトリが赤ん坊を運んでくるわけではない. 国の面積や人口といった「国の規模」が両方を押し上げているだけだ.
交絡因子は気温のような物理量だけでなく, 国の規模のような構造量でも起こる.

次に, $x$が$y$を動かしているのではなく, $y$が$x$を動かしているケース. \textbf{逆因果}と呼ぶ.
たとえば「医療支出が多い地域ほど健康指標が悪い」というデータが出たとしよう.
医療支出が悪さをしているのではなく, 健康状態が悪い地域だからこそ医療支出が増えている, という方向の因果かもしれない.
矢印の向きを取り違えると, 結論はそのまま逆転する.

そして, そもそも何の関係もない\textbf{偶然の一致}.
無関係な2つの量を大量に集めて相関を計算すれば, そのうちのいくつかは偶然強く相関する.
「米国のチーズ消費量とベッドシーツでの絞死者数」のような, ふざけた相関集が世間にあるくらいだ.
たくさん試せば偽陽性は当然出る, という前節のp-ハッキングと同じ構造の話である.

強い相関を見たときに自動的に問うべきは, 裏で動いている第3変数はないか, 因果の向きはどちらか, 似たような量を大量に試した結果ではないか.
このどれにも即答できないなら,「強い相関がある」というだけの結論で立ち止まらないほうがよい.

\subsubsection*{データの外側に踏み込むとき}

8章で広告費と売上の回帰直線$y \approx 119 + 2.3x$を引いた.
「来月の広告費を50万円にしたら売上はおよそ234万円」と予測した.
50万円はデータの範囲（20万円から80万円）に収まっているから, この予測には根拠がある.

では, 広告費200万円ならどうか.
式に代入すれば579万円だ. 計算は問題なく動く.
しかし200万円はデータの範囲をはるかに超えた値である.
広告費20万円から80万円の範囲で観測された直線関係が, 200万円まで続いている保証はどこにもない.
広告には効果の頭打ちがある場面が多いし, 別のメカニズムが働き始める場合もある.

これが\textbf{外挿}の問題だ.
データの範囲外に直線を延ばして予測すると,「直線が続いている」という根拠のない仮定を黙って置くことになる.
電卓は何も警告しない. 警告するのは使い手だ.

切片の解釈も外挿と地続きの話だ.
$a \approx 119$は, 式の上では「広告費がゼロのときの売上119万円」となる.
しかし広告費ゼロの月は実データにない.
この119万円は, 直線の位置を決めるために計算上現れた数値である. 現実の予測値ではなく, 直線をy軸まで延ばしたときの幾何学的な交点だと考えるのが安全だ.
報告するときも, 切片そのものを「広告ゼロでも売れる金額」のように紹介してはいけない.「観測した広告費20万円から80万円の範囲で, 傾き2.3, 切片119の直線がよく当てはまる」という言い方に留める.

「数字が出る」と「現実に当てはめてよい」は別だ.
モデルが描く直線は, 観測した範囲で関係を要約する道具であって, その外側まで現実を保証するものではない.

ここまでが, 計算は正しいのに結論が間違う典型だった.
最後は, 数字を扱うこと自体に伴う責任の話に視野を広げて, 本書を閉じる.


\subsection{データを扱う責任と, この本の終わりに}

ここまでは技術と解釈の話だった.
最後に1つ, 技術の続きとしての話をしておく. 倫理の話だ.

倫理という言葉は, 統計の教科書の中ではいかにも浮く.
ところが本章で並べた歪みは, 個人がうっかり犯すだけで終わらず, 社会の規模で誰かに不利益を与える形にもなりうる.
そこに歯止めをかけるのが, データを扱う側の責任の話だ.

\subsubsection*{設計の逸脱が社会保障に直結した例}

日本では2018年に, 厚生労働省の毎月勤労統計をめぐる問題が発覚した.
毎月勤労統計は雇用保険や労災保険の給付額を決める基礎資料の1つであり, 法律上は従業員500人以上の大規模事業所は全数調査することになっていた.
ところが東京都分について, 2004年から長年, 全数調査ではなく一部の抽出調査が行われていた.
さらに, 抽出に伴って必要な復元処理（拡大処理）が抜けていた.

ここで「復元処理」を補っておこう.
全数調査するはずだった事業所のうち一部しか抽出していないなら, 抽出された一部の値をそのまま全体の代表として使うわけにはいかない.
抽出した事業所の合計を,「もし全数だったら」の規模に拡大しなおす作業が復元処理だ.
たとえば3分の1だけ抽出した分を3倍に直す, というイメージでよい.
東京都の大規模事業所には平均賃金が高い企業が多い. 復元せずに全国集計に混ぜれば, その「賃金の高い部分」が本来より小さく取り込まれ, 全国平均が低めに計上される.

この狂いが雇用保険や労災保険の給付額に響いた.
影響を受けた人は延べ約2{,}000万人, 過少支給の総額は約500億円超と公表されている.
追加給付のための事務費は別途数百億円規模で発生した.

「全数調査するはずの場所で抽出をしていた」という設計の逸脱が, 社会保障の配分という実体に直結した事例だ.
9.2で扱ったサンプリングバイアスが, 個人の研究の話ではなく公的な仕組みの中で起きていたわけだ.

\subsubsection*{最低限の3点：同意と目的, 保護, 透明性}

データを扱うときの規範は, 国や分野によって細部が違うが, 大枠は3点に集約できる.

最初は, 集めるときに目的を明示し, 同意を得ること.
何のためのデータなのか, どこまで使われるのかを, 提供する側が理解できる形で伝える.
あとから目的外に流用しない.

次に, 個人が特定できる情報を保護すること.
氏名や住所そのものを直接保管するかどうかの話だけではない.
複数の項目を組み合わせると個人が特定できてしまう場合があるので, 公開する集計の単位や粒度には注意がいる.

そして, 報告の透明性. 手順, 前処理, 除外したデータの基準を, 第三者が検証できる形で公開する.
9.2で扱った確証バイアスやp-ハッキングは, 透明性の欠如の上にこそ繁茂する.
誰がどの段階で何を選んだかが見えていれば, 外側からの検証によって歪みは正されうる.

3点はバラバラのルールではない.
9.2で見た上流・中流・下流の歪みを, 社会の側で抑える仕組みだと考えればよい.
集め方を歪ませない（上流対策）, 個人を巻き添えにしない（中流対策）, 解釈の自由を残しすぎない（下流対策）.
本章で習った歪みの地図と, ここでの3点は, 同じ地形を別の角度から眺めた地図である.

\subsubsection*{ここまでに何ができるようになったか}

本書を1章から振り返ると, ずいぶん遠くまで来た.
2章でデータの大きさと散らばりを言葉にし, 3章で形を見る道具を手にした.
4章で「取り直したらばらつく」という事情を引き受け, CLTを通じて正規分布が現れる理屈を理解し, SEで推定値の幅を測れるようになった.
5章ではその幅を信頼区間として運用する技を, 6章では「偶然では説明しにくいかどうか」を確かめる検定の枠組みを得た.
7章で2変数の関係の方向と強さを相関係数で測り, 8章で関係を直線で表し, $R^2$と残差プロットで自分のモデルを疑う習慣をつくった.
そして本章で, 手法の外側にある歪みを地図にした.

ここまでに獲得した能力を一言で表すなら,「数字が出る」から「数字を疑える」までの距離だ.
学び始めの頃は, 計算結果を出せるようになることが目標だったかもしれない.
今となっては, それは出発点にすぎなかったとわかる.
出した数字の前提は何か. その前提は満たされているか. 集計の単位は妥当か. 第3変数を見落としていないか. データの範囲外まで延ばして語っていないか.
数字を出した後にこれだけ問えるようになっていれば, 本書の役目はおおむね果たされた.

\subsubsection*{次の一歩}

本書は統計の骨格に絞った. 補足概念として各章の末尾に名前だけ挙げた話題は, それぞれ独立した教科書1冊分の世界が広がっている.
多重比較の制御, ブートストラップによる区間推定, 重回帰と多重共線性, 時系列分析の独立性の扱い方, 事前登録と再現性の手順, データガバナンスとプライバシ.

順序の目安だけ示しておく.
まず, 本書で繰り返し触れたp値の正しい運用に直結する\textbf{多重比較の制御}と\textbf{事前登録}が, 6章の延長として一番近い.
次に, 8章の直線回帰を拡張する\textbf{重回帰と多重共線性}が, 関係を測る道具を厚くしてくれる.
分布や区間推定をもう一段柔らかくしたいなら\textbf{ブートストラップ}, 時間軸が絡むデータを扱うなら\textbf{時系列分析}, 業務でデータを集める側に回るなら\textbf{データガバナンスとプライバシ}.
本書で身につけた骨格の上に, 必要を感じた順に積み上げていけばよい.

最後に1つだけ提案する.
読者自身の手元のデータで, 本書の手順を最初から最後まで一度回してみてほしい.
要約し, 形を見て, 幅を測り, 差を吟味し, 関係を見て, モデルにして, 出した結論を疑う.
その過程で, 本書のどこかが「もっと知りたい」「ここはなぜそうなのか」と引っかかるはずだ.
その引っかかりこそが, 統計学の次の章への入口になる.

ここから先は, あなたが自分のデータと向き合う番である.
